%==========================================================================
%  TRES0312 Fysikk — Foiler-mal
%
%  MØrk/lys modus: kommenter ut én av de to linjene nedenfor.
%==========================================================================
\newif\ifdarkmode
\darkmodetrue      % ← mørk modus
%\darkmodefalse   % ← lys modus  (fjern % foran og legg % foran linjen over)

\documentclass[aspectratio=169, 12pt]{beamer}

\usepackage{amd-beamer}
\usetikzlibrary{overlay-beamer-styles}
\usetikzlibrary{calc}

%--------------------------------------------------------------------------
%  DOKUMENTINFO
%--------------------------------------------------------------------------
\title{Trykk i væsker og Pascals lov}
\subtitle{TRES0312 --- kapittel 7.2 og 7.3}
\author{Ben David Normann}
\date{\today}

%--------------------------------------------------------------------------
%  SEKSJONSKONTEKST
%--------------------------------------------------------------------------
\setcounter{section}{7}   % Kapittel 7: "Fluiddynamikk"

\begin{document}

%--------------------------------------------------------------------------
%  SLIDE 1 — TITTELSIDE
%--------------------------------------------------------------------------
\begin{frame}[plain]
  \titlepage
  \dekkerdelkap{kapittel 7.2 og 7.3}
  \vspace{0.3cm}
  \begin{center}
  {\bfseries\color{repcol} Første time: trykk i væsker. Etter pausen: Pascals lov og hydrauliske systemer.}
  \end{center}
\end{frame}

%--------------------------------------------------------------------------
%  SLIDE 2 — Tre kjappe fra forrige gang
%--------------------------------------------------------------------------
\begin{frame}{}

  \repetisjon{Tre kjappe fra forrige gang}{%
    \begin{enumerate}[1.]
      \item Hva er definisjonen av massetetthet? Hva er massen til 1 L vann?
      \item Hva er definisjonen av trykk? Er trykk en vektor?
      \item Omtrent hvor stort er lufttrykket, målt i pascal?
    \end{enumerate}
  }

\end{frame}

%==========================================================================
%  DEL 1: TRYKK I VÆSKER (7.2)
%==========================================================================


%--------------------------------------------------------------------------
%  SLIDE 3 — Til ettertanke: to punkter på samme dybde
%--------------------------------------------------------------------------
\begin{frame}{Horisontalt væskeelement}

  \tanke{}{%
  To punkter ligger på samme dybde i et basseng der vannet står helt i ro. Kan trykket være større i det ene punktet enn i det andre?
  }
  \pause
  \svar{
  Nei. Hvis trykket var større i det ene punktet, ville vannet begynt å strømme mot punktet med lavere trykk. Men vannet står i ro. Altså må trykket være likt.
  }

\end{frame}

%--------------------------------------------------------------------------
%  SLIDE 4 — Definisjon: trykk i horisontalt væskeelement
%--------------------------------------------------------------------------
\begin{frame}{Horisontalt væskeelement}

  \defnat{7.2.1}{Trykk i horisontalt væskeelement}{
  I en væske i ro er trykket likt i alle punkter på samme horisontale nivå:
  \[
    p_1 = p_2
  \]
  Punktene må ligge i den \textit{samme sammenhengende} væsken.
  }

  \pause
  \merk{Vi holder oss til væsker i ro}{%
  Når en væske strømmer, kan trykket variere langs samme nivå. Der et rør er smalt, strømmer væsken fortere, og trykket er lavere. Det er ikke pensum i dette faget.
  }

\end{frame}

%--------------------------------------------------------------------------
%  SLIDE 5 — Aktivitet: trykket i et akvarium
%--------------------------------------------------------------------------
%  Fasit: a) p_B = p_A = 1.05e5 Pa. b) p_C = p_A = 1.05e5 Pa (samme nivå, samme
%         sammenhengende vann, selv om det er mindre vann rett over C).
%         c) p_D > p_A (mer vann over D).
\begin{frame}{Aktivitet}

  \begin{columns}[T]
    \begin{column}{0.50\textwidth}
      \aktivitet{Trykket i et akvarium}{%
      Punktene A, B og C ligger $0.40\,\mathrm{m}$ under vannflaten. C ligger inne i en hule under en stein. Trykket i A er $1.05\cdot 10^5\,\mathrm{Pa}$.
      \begin{enumerate}[a)]
        \item Hva er trykket i B?
        \item Hva er trykket i C?
        \item Er trykket i D større eller mindre enn i A?
      \end{enumerate}
      }
    \end{column}
    \begin{column}{0.48\textwidth}
      \centering
      \vspace{6mm}
      \begin{tikzpicture}[scale=0.85, font=\small]
        \fill[oppgcol!45!pagecolor] (0,0) rectangle (7.0,3.0);
        \fill[gray!55, rounded corners=4pt] (5.0,0) -- (5.0,1.4) -- (3.8,1.4) -- (3.8,1.9) -- (4.5,2.2) -- (6.2,2.25) -- (6.6,1.8) -- (6.6,0) -- cycle;
        \draw[very thick, textcolor] (0,3.5) -- (0,0) -- (7.0,0) -- (7.0,3.5);
        \draw[oppgcol!70!white, thick] (0,3.0) -- (7.0,3.0);
        \draw[dashed, textcolor!60] (0,1.0) -- (7.0,1.0);
        \foreach \x/\n in {1.2/A, 2.8/B, 4.4/C}{
          \fill[textcolor] (\x,1.0) circle (2pt);
          \node[below] at (\x,1.0) {\n};
        }
        \fill[textcolor] (2.0,0.35) circle (2pt);
        \node[right] at (2.0,0.35) {D};
      \end{tikzpicture}
    \end{column}
  \end{columns}

\end{frame}

%--------------------------------------------------------------------------
%  SLIDE 6 — Definisjon: trykk i vertikalt væskeelement
%--------------------------------------------------------------------------
\begin{frame}{Vertikalt væskeelement}

  \defnat{7.2.2}{Trykk i vertikalt væskeelement}{
  Trykket i dybden $h$ i en væske i ro er
  \[
    p(h) = p_0 + \rho\mathrm{g}h
  \]
  der $p_0$ er trykket ved overflaten og $\rho$ er tettheten til væsken.
  }

  \pause
  \merk{Hydrostatisk trykk}{%
  Leddet $\rho\mathrm{g}h$ kalles det \textit{hydrostatiske trykket}. Det skyldes vekten av væsken over punktet. Mellom to punkter med høydeforskjell $h$ er trykkforskjellen $p_2 - p_1 = \rho\mathrm{g}h$.
  }

\end{frame}

%--------------------------------------------------------------------------
%  SLIDE 7 — Til ettertanke: tre beholdere
%--------------------------------------------------------------------------
\begin{frame}{Vertikalt væskeelement}

  \begin{columns}[c]
    \begin{column}{0.58\textwidth}
      \tanke{}{%
      Tre beholdere med ulik form fylles med vann til samme høyde. Er trykket i bunnen størst i beholderen med mest vann?
      }
      \pause
      \svar{
      Nei. Trykket i bunnen er $p_0 + \rho\mathrm{g}h$ i alle tre. Det avhenger bare av dybden, ikke av formen eller hvor mye vann det er.
      }
    \end{column}
    \begin{column}{0.40\textwidth}
      \centering
      \begin{tikzpicture}[scale=0.8]
        % Smal sylinder
        \fill[oppgcol!45!pagecolor] (0,0) rectangle (0.7,2.2);
        \draw[very thick, textcolor] (0,2.6) -- (0,0) -- (0.7,0) -- (0.7,2.6);
        % Vid bolle
        \fill[oppgcol!45!pagecolor] (1.3,0) -- (1.0,2.2) -- (3.0,2.2) -- (2.7,0) -- cycle;
        \draw[very thick, textcolor] (0.95,2.6) -- (1.3,0) -- (2.7,0) -- (3.05,2.6);
        % Kjegle med spissen opp
        \fill[oppgcol!45!pagecolor] (3.4,0) -- (4.15,2.2) -- (4.45,2.2) -- (5.2,0) -- cycle;
        \draw[very thick, textcolor] (3.4,0) -- (5.2,0) -- (4.36,2.6) (3.4,0) -- (4.24,2.6);
        \draw[dashed, textcolor!60] (-0.2,2.2) -- (5.4,2.2);
      \end{tikzpicture}
    \end{column}
  \end{columns}

\end{frame}

%--------------------------------------------------------------------------
%  SLIDE 8 — Aktivitet: dykker
%--------------------------------------------------------------------------
%  Fasit: p = 101325 + 1000*9.81*h.
%         h = 3.0 m: 1.31e5 Pa ≈ 1.3e5 Pa.
%         h = 8.0 m: 1.80e5 Pa ≈ 1.8e5 Pa.
%         h = 20 m:  2.98e5 Pa ≈ 3.0e5 Pa.
%         Ca. 1 atm ekstra per 10 m.
\begin{frame}{Aktivitet}

  \aktivitet{Trykket på en dykker}{%
  En dykker er under vann ($\rho = 1000\,\mathrm{kg/m^3}$). Lufttrykket ved overflaten er $p_0 = 1.013\cdot 10^5\,\mathrm{Pa}$.
  \begin{enumerate}[a)]
    \item Finn det totale trykket på dykkeren på dybdene $3.0\,\mathrm{m}$, $8.0\,\mathrm{m}$ og $20\,\mathrm{m}$.
    \item Omtrent hvor mye øker trykket for hver tiende meter nedover?
  \end{enumerate}
  }

\end{frame}

%--------------------------------------------------------------------------
%  SLIDE 9 — Aktivitet: rør heises opp av vannet
%--------------------------------------------------------------------------
%  Fasit: a) h = p_0/(rho g) = 1.013e5/(1000*9.81) = 10.3 m.
%         b) Tilnærmet vakuum (litt vanndamp, ca. 2 kPa), så i praksis ca. 10.1 m.
\begin{frame}{Aktivitet}

  \begin{columns}[T]
    \begin{column}{0.62\textwidth}
      \aktivitet{Et langt rør heises opp av vannet}{%
      Et $15\,\mathrm{m}$ langt rør er tett i den ene enden. Det ligger helt fylt med vann i sjøen. En kran heiser røret loddrett opp, med den tette enden først. Den åpne enden er hele tiden under vannflaten.
      \begin{enumerate}[a)]
        \item Hvor høy blir vannsøylen i røret, målt fra vannflaten utenfor?
        \item Over vannsøylen ser det ut som om det er luft. Hva er det egentlig der?
      \end{enumerate}
      }
    \end{column}
    \begin{column}{0.36\textwidth}
      \centering
      \begin{tikzpicture}[scale=0.21, font=\footnotesize, >={Stealth[length=5pt,width=4pt]}]
        \fill[oppgcol!45!pagecolor] (-4,-3) rectangle (5,0);
        \draw[oppgcol!70!white, thick] (-4,0) -- (5,0);
        \fill[oppgcol!65!pagecolor] (0,-2) rectangle (0.9,10.3);
        \draw[very thick, textcolor] (0,-2) -- (0,15) -- (0.9,15) -- (0.9,-2);
        \draw[thick, textcolor!70] (0.45,15) -- (0.45,17);
        \fill[textcolor!70] (0.45,17) circle (0.2);
        \node[font=\small\bfseries] at (0.45,12.6) {?};
        \draw[<->, thin, textcolor!70] (2.2,0) -- (2.2,10.3) node[midway, right] {$h = {?}$};
        \draw[dashed, textcolor!50] (0.9,10.3) -- (2.6,10.3);
      \end{tikzpicture}
    \end{column}
  \end{columns}

\end{frame}

\begin{frame}{Vertikalt væskeelement}

  \tanke{}{%
  Hva er det som holder vannsøylen i røret oppe?
  }
  \pause
  \svar{
  Lufttrykket på vannflaten \textit{utenfor} røret. Inne i røret, i høyde med vannflaten utenfor, er trykket $p_0$. Over vannsøylen er det (nesten) ikke noe trykk. Da kan lufttrykket holde oppe en vannsøyle på $h = p_0/(\rho\mathrm{g}) \approx 10\,\mathrm{m}$.
  }

\end{frame}

%--------------------------------------------------------------------------
%  SLIDE 10 — Oppsummering: trykk i væsker
%--------------------------------------------------------------------------
\begin{frame}{}

  \oppsum{Trykk i væsker}{%
  \begin{enumerate}[1.]
    \item[] I en væske i ro er trykket likt på samme nivå: $p_1 = p_2$ (samme sammenhengende væske).
    \item[] Trykket øker med dybden: $p(h) = p_0 + \rho\mathrm{g}h$.
    \item[] Trykket avhenger bare av dybden, ikke av formen på beholderen.
    \item[] I vann øker trykket med ca. 1 atm for hver tiende meter.
  \end{enumerate}
  }

\end{frame}

%--------------------------------------------------------------------------
%  SLIDE 11 — PAUSE
%--------------------------------------------------------------------------
\begin{frame}
\begin{center}
  \Huge{Pause}\\[8pt]
  {\large Etter pausen: Pascals lov og hydrauliske systemer}
\end{center}
\end{frame}

%==========================================================================
%  DEL 2: PASCALS LOV OG HYDRAULISKE SYSTEMER (7.3)
%==========================================================================


%--------------------------------------------------------------------------
%  SLIDE 12 — Til ettertanke: tannkremtuben
%--------------------------------------------------------------------------
\begin{frame}{Pascals lov}

  \tanke{}{%
  Du klemmer på midten av en tannkremtube. Hvorfor kommer tannkremen ut i enden, og ikke bare der du klemmer?
  }
  \pause
  \svar{
  Trykkøkningen fra fingrene forplanter seg gjennom hele tannkremen, helt ut til åpningen.
  }

\end{frame}

%--------------------------------------------------------------------------
%  SLIDE 13 — Pascals lov
%--------------------------------------------------------------------------
\begin{frame}{Pascals lov}

  \begin{itemize}
    \item Øker vi trykket på en innestengt væske i ro med $\Delta p$, øker trykket \textit{overalt} i væsken med den samme $\Delta p$.
    \item<2-> Det totale trykket varierer fortsatt med dybden. I dybden $h$ under overflaten er trykket
  \end{itemize}
  \onslide<2->{\graabox{$\displaystyle p(h) = p_0 + \rho\mathrm{g}h + \Delta p$}}
  \onslide<3->{\merk{Pascals lov er grunnlaget for hydrauliske systemer}{%
  Der kan en liten kraft forsterkes ved hjelp av et stort areal. Det ser vi om litt.
  }}

\end{frame}

%--------------------------------------------------------------------------
%  SLIDE 14 — Aktivitet: lodd på et stempel
%--------------------------------------------------------------------------
%  Fasit: a) Δp = m g / A = 30*9.81/0.020 = 1.47e4 Pa ≈ 1.5e4 Pa.
%         b) Samme: 1.5e4 Pa (Pascals lov).
%         c) p = 1.013e5 + 1.47e4 + 1000*9.81*0.40 = 1.20e5 Pa ≈ 1.2e5 Pa.
\begin{frame}{Aktivitet}

  \begin{columns}[T]
    \begin{column}{0.72\textwidth}
      \aktivitet{Lodd på et stempel}{%
      En beholder med vann er lukket med et tett stempel med areal $A = 0.020\,\mathrm{m^2}$. Over stempelet er det luft med trykket $p_0 = 1.013\cdot 10^5\,\mathrm{Pa}$. Vi legger et lodd på $30\,\mathrm{kg}$ på stempelet.
      \begin{enumerate}[a)]
        \item Hvor mye øker trykket like under stempelet?
        \item Hvor mye øker trykket i bunnen, $0.40\,\mathrm{m}$ lenger ned?
        \item Hva er det totale trykket i bunnen?
      \end{enumerate}
      }
    \end{column}
    \begin{column}{0.26\textwidth}
      \centering
      \vspace{6mm}
      \begin{tikzpicture}[scale=0.75, font=\small]
        \fill[oppgcol!45!pagecolor] (0,0) rectangle (2.0,2.5);
        \fill[gray!55] (0,2.5) rectangle (2.0,2.75);
        \fill[gray!80] (0.6,2.75) rectangle (1.4,3.35);
        \node[font=\footnotesize, white] at (1.0,3.05) {$30\,\mathrm{kg}$};
        \draw[very thick, textcolor] (0,3.6) -- (0,0) -- (2.0,0) -- (2.0,3.6);
        \draw[<->, thin, textcolor!70] (2.4,2.5) -- (2.4,0) node[midway, right, font=\footnotesize] {$0.40\,\mathrm{m}$};
      \end{tikzpicture}
    \end{column}
  \end{columns}

\end{frame}

%--------------------------------------------------------------------------
%  SLIDE 15 — Til ettertanke: løfte en bil
%--------------------------------------------------------------------------
\begin{frame}{Hydrauliske systemer}

  \tanke{}{%
  Hvordan kan du løfte en bil på over et tonn med en jekk, bare ved å presse lett på et lite stempel?
  }
  \pause
  \svar{
  Trykkøkningen $\Delta p$ er den samme ved begge stemplene. Men kraft er trykk ganget med areal, $F = \Delta p\cdot A$. Et stort stempel får derfor en stor kraft.
  }

\end{frame}

%--------------------------------------------------------------------------
%  SLIDE 16 — Definisjon: Pascals lov i hydrauliske systemer
%--------------------------------------------------------------------------
\begin{frame}{Hydrauliske systemer}

  \begin{columns}[c]
    \begin{column}{0.58\textwidth}
      \defnat{7.3.1}{Pascals lov i hydrauliske systemer}{
      For to stempler koblet sammen gjennom samme væske er
      \[
        \frac{F_1}{A_1} = \frac{F_2}{A_2}
        \quad\Longrightarrow\quad
        F_2 = F_1\cdot\frac{A_2}{A_1}
      \]
      }
    \end{column}
    \begin{column}{0.40\textwidth}
      \centering
      \begin{tikzpicture}[scale=0.85, font=\small, >={Stealth[length=7pt,width=5pt]}]
        \fill[oppgcol!45!pagecolor] (0,0) -- (4.6,0) -- (4.6,1.2) -- (1.6,1.2) -- (1.6,0.4) -- (0.6,0.4) -- (0.6,1.2) -- (0,1.2) -- cycle;
        \fill[gray!55] (0,1.2) rectangle (0.6,1.4);
        \fill[gray!55] (1.6,1.2) rectangle (4.6,1.4);
        \draw[very thick, textcolor] (0,2.4) -- (0,0) -- (4.6,0) -- (4.6,2.4);
        \draw[very thick, textcolor] (0.6,2.4) -- (0.6,0.4) -- (1.6,0.4) -- (1.6,2.4);
        \draw[->, line width=2.5pt, red!70!white] (0.3,1.4) -- (0.3,0.55);
        \node at (0.3,1.9) {$\vec{F}_1$};
        \draw[->, line width=2.5pt, defscol] (3.1,1.2) -- (3.1,2.9) node[right=2pt] {$\vec{F}_2$};
        \node[font=\footnotesize] at (0.3,2.65) {$A_1$};
        \node[font=\footnotesize] at (4.1,1.75) {$A_2$};
      \end{tikzpicture}
    \end{column}
  \end{columns}

\end{frame}

%--------------------------------------------------------------------------
%  SLIDE 17 — Aktivitet: hydraulisk jekk
%--------------------------------------------------------------------------
%  Fasit: F_2 = 60 N * 160/4.0 = 2.4e3 N (40 ganger større).
\begin{frame}{Aktivitet}

  \aktivitet{Hydraulisk jekk}{%
  En hydraulisk jekk har et lite stempel med areal $A_1 = 4.0\,\mathrm{cm^2}$ og et stort stempel med areal $A_2 = 160\,\mathrm{cm^2}$. Du presser på det lille stempelet med kraften $F_1 = 60\,\mathrm{N}$.
  \begin{enumerate}[a)]
    \item Finn kraften $F_2$ på det store stempelet.
    \item Hvor mange ganger større er $F_2$ enn $F_1$?
  \end{enumerate}
  }

\end{frame}

%--------------------------------------------------------------------------
%  SLIDE 18 — Til ettertanke: får vi noe gratis?
%--------------------------------------------------------------------------
\begin{frame}{Hydrauliske systemer}

  \tanke{}{%
  Jekken gir oss 40 ganger så stor kraft. Får vi da noe «gratis»?
  }
  \pause
  \svar{
  Nei. Væsken som presses ut av den lille sylinderen, må inn i den store. Det lille stempelet må derfor presses 40 ganger så langt som det store løftes. Arbeidet $W = F\cdot s$ er det samme på begge sider: energien er bevart.
  }

\end{frame}

%--------------------------------------------------------------------------
%  SLIDE 19 — Oppsummering: Pascals lov og hydrauliske systemer
%--------------------------------------------------------------------------
\begin{frame}{}

  \oppsum{Pascals lov og hydrauliske systemer}{%
  \begin{enumerate}[1.]
    \item[] En trykkøkning $\Delta p$ forplanter seg uforminsket til hele væsken: $p(h) = p_0 + \rho\mathrm{g}h + \Delta p$.
    \item[] Samme trykkøkning ved begge stemplene: $F_1/A_1 = F_2/A_2$, så $F_2 = F_1\cdot A_2/A_1$.
    \item[] Men det lille stempelet må flyttes tilsvarende lenger. Arbeidet er det samme.
  \end{enumerate}
  }

\end{frame}

%--------------------------------------------------------------------------
%  SLIDE 20 — Neste gang
%--------------------------------------------------------------------------
\begin{frame}
\begin{center}
  \Huge{Neste gang: oppdrift og Arkimedes' lov}
  \vspace{6mm}
\end{center}
\end{frame}

\end{document}
